\documentclass{article}
\usepackage{iclr2026_conference,times}

%%%%% Math definitions for spectral_partition_router_v2_1 %%%%%

\usepackage{amsmath,amsfonts,bm}

\newcommand{\newterm}[1]{{\bf #1}}
\def\figref#1{figure~\ref{#1}}
\def\Figref#1{Figure~\ref{#1}}
\def\tabref#1{table~\ref{#1}}
\def\Tabref#1{Table~\ref{#1}}
\def\secref#1{section~\ref{#1}}
\def\Secref#1{Section~\ref{#1}}
\def\eqref#1{equation~\ref{#1}}
\def\Eqref#1{Equation~\ref{#1}}

\newcommand{\R}{\mathbb{R}}
\newcommand{\E}{\mathbb{E}}
\newcommand{\Prob}{\mathbb{P}}
\newcommand{\Sphere}{\mathbb{S}}
\newcommand{\Cov}{\mathrm{Cov}}
\newcommand{\dist}{\mathrm{dist}}
\newcommand{\calD}{\mathcal{D}}
\newcommand{\calE}{\mathcal{E}}
\newcommand{\calN}{\mathcal{N}}
\newcommand{\calT}{\mathcal{T}}
\newcommand{\calC}{\mathcal{C}}
\newcommand{\calG}{\mathcal{G}}
\newcommand{\calP}{\mathcal{P}}
\newcommand{\norm}[1]{\left\lVert #1 \right\rVert}
\newcommand{\inner}[2]{\left\langle #1,#2 \right\rangle}


\usepackage{hyperref}
\usepackage{url}
\usepackage{booktabs}
\usepackage{amsmath}
\usepackage{amssymb}
\usepackage{amsthm}
\usepackage{graphicx}
\usepackage{microtype}
\usepackage{placeins}
\usepackage{flafter}
\hypersetup{hidelinks}

\setlength{\textfloatsep}{8pt plus 2pt minus 2pt}
\setlength{\floatsep}{8pt plus 2pt minus 2pt}
\setlength{\intextsep}{8pt plus 2pt minus 2pt}

\newtheorem{definition}{Definition}
\newtheorem{assumption}{Assumption}
\newtheorem{proposition}{Proposition}
\newtheorem{theorem}{Theorem}
\newtheorem{remark}{Remark}

\makeatletter
\renewenvironment{proof}[1][\proofname]{\par\noindent\textit{#1.}\ }{\hfill\qedsymbol\par}
\makeatother

\title{Geometry-Aware Spectral Partition Routing for Uniform Feature Dispatch}

\author{Anonymous Authors}

%\iclrfinalcopy
\begin{document}

\maketitle

\begin{abstract}
We study a controlled routing problem for mixture-of-experts models: when
route-relevant features are spectrally uniform, can we build a non-cheating
router whose dispatch is uniform, feature-aware, and stable during training?
Prior equal-mass constructions solve only the first part: quantile thresholds can
make expert cells have probability $1/E$, but equal mass alone does not imply
that the cells align with latent feature structure or survive optimization.  We
therefore formulate \emph{good feature dispatch} as three separable properties:
equal mass, feature coherence, and margin-based preservation.  We propose
Geometry-Aware Spectral Partition Routing, which keeps recursive quantile splits
for mass control, but chooses split directions using either unlabeled locality or
expert-weight geometry, and adds margin-aware preservation rules for training.
The paper gives the exact equal-mass proof, a discrete-feature tie boundary, a
coherence guarantee under projected separation, and a preservation theorem under
bounded representation drift.  The experiment section is a protocol only: it
specifies static, feature-structure, expert-guided, and training-preservation
audits to be filled after execution.
\end{abstract}

\section{Introduction}

Mixture-of-experts (MoE) models route each token representation to a small set of
experts \citep{shazeer2017outrageously,fedus2022switch}.  In a top-1 router,
small geometric differences decide the selected expert.  Therefore, if the
research question is whether uniformly represented features can be distributed to
experts, the first object to define is not the load-balancing loss, but the
feature-space partition induced by the router.

The simplest useful setting is controlled: route-relevant features have been
chosen, centered, and spectrally normalized so that no dominant covariance
direction remains.  This is not a claim about real language-model hidden states;
it is the clean case one should solve before studying nonuniform spectra,
residual common components, or sequence-position mixing.  The remaining question
is constructive: what router should we use, and what properties should it have?

A purely equal-mass answer is too weak.  Recursive quantile splitting can build a
fixed feature router $g(x)$ whose cells $C_e=\{x:g(x)=e\}$ satisfy
$\Prob(C_e)=1/E$.  This solves count-level dispatch without using labels or
batch-level capacity forcing.  However, experiments from the previous iteration
showed that equal-mass cells can have weak feature agreement, and that naive
recalibration can destroy route stability.  Thus a ``good'' router must satisfy
three properties, not one: it must dispatch equal probability mass, cut features
in a geometrically meaningful way when structure exists, and preserve the
assignment under training drift.

This paper rewrites the problem around those three properties.  We propose
\newterm{Geometry-Aware Spectral Partition Routing} (GA-SPR).  GA-SPR keeps the
quantile mechanism because it is the clean source of equal mass.  It then adds a
geometry layer: split directions are chosen from unlabeled feature locality or
from expert weights.  Expert weights are allowed because they are model-internal
information, not data labels.  In particular, an expert matrix supplies a
candidate spectral prototype through its principal input singular direction,
matching the expert-proxy view of spectral router designs such as MPI
\citep{routermpi2026}.  Finally, GA-SPR treats training preservation as a margin
problem: updates should be accepted only when they do not induce excessive route
flips near partition boundaries.

\paragraph{Contributions.}
\begin{enumerate}
\item We refine uniform feature dispatch into three measurable targets: equal
mass, feature coherence, and margin-based preservation.  This clarifies what is
solved by quantiles and what requires additional structure.
\item We define GA-SPR, a family of non-cheating routers that combine recursive
quantile mass control with locality-aware or expert-weight-guided split
directions, plus margin-aware training updates.
\item We prove the mass, coherence, and preservation claims under explicit mild
assumptions, and specify an experiment protocol that keeps static dispatch,
feature quality, expert-weight alignment, and training dynamics separate.
\end{enumerate}

\section{Problem Definition}

\subsection{Controlled feature setting}

Let $X\in\R^d$ be the route-relevant feature.  In a sequence model this should be
the hidden state at the position used for routing, not a pool of all positions.
We assume an unlabeled calibration set
\begin{equation}
  \calD_{\mathrm{cal}}=\{x_i\}_{i=1}^{n_{\mathrm{cal}}}
\end{equation}
from the same feature population.  Calibration may use feature statistics and
model-internal expert weights, but not feature labels or task targets.

The representation is spectrally normalized by
\begin{equation}
  \tilde x = B(x-\mu),\qquad \Cov(\tilde X)\approx \sigma^2 I.
\end{equation}
Spectral normalization removes a dominant covariance direction as a confound.  It
does not by itself imply that the feature distribution has useful local
structure.

\subsection{Three targets for good dispatch}

A router is a deterministic map $g:\R^d\rightarrow\{1,\ldots,E\}$ with cells
$C_e=\{x:g(x)=e\}$.

\begin{definition}[Equal-mass dispatch]
The router has equal-mass dispatch if
\begin{equation}
  \Prob(X\in C_e)=\frac{1}{E},\qquad e=1,\ldots,E.
\end{equation}
\end{definition}

Equal mass is the base requirement.  It is not sufficient for a good router.

\begin{definition}[Feature coherence]
Suppose the data distribution has latent feature identity $F$ for analysis only.
A router is feature-coherent when $g(X)$ is informative about $F$, measured by
route-feature mutual information, purity, or low within-cell variance.  The
method may not use $F$ during calibration or training.
\end{definition}

When no latent structure exists, feature coherence is not identifiable from
isotropy alone.  Therefore coherence claims require an additional assumption such
as clusterable features or expert weights aligned with feature directions.

\begin{definition}[Training preservation]
Let $g_0$ be the calibrated route and $g_t$ be the route after training step $t$.
The partition is preserved if the route flip probability
\begin{equation}
  \Prob(g_t(X_t)\ne g_0(X_0))
\end{equation}
is small, and if the lower tail of the routing margin remains away from zero.
\end{definition}

\subsection{Non-cheating rule}

A router is non-cheating if, after calibration, routing a token uses only the
current feature, stored calibration statistics, and optional model-internal
expert weights.  It may not use feature labels, task targets, evaluation labels,
or current-batch expert counts.  Batch-level balanced assignment is therefore a
control, not the method.

\section{Geometry-Aware Spectral Partition Routing}

GA-SPR is a recursive tree router.  Each internal node owns a feature
subpopulation and a subset of experts.  Each leaf owns one expert.

\subsection{Mass layer: recursive quantile split}

At node $\calN$, let $\calE_{\calN}$ be the current expert set and split it into
$\calE_L$ and $\calE_R$.  The target left mass is
\begin{equation}
  \rho_{\calN}=\frac{|\calE_L|}{|\calE_{\calN}|}.
\end{equation}
Given a direction $q_{\calN}\in\Sphere^{d-1}$, define
\begin{equation}
  T_{\calN}=q_{\calN}^{\top}\tilde X,\qquad \tilde X\sim P_{\calN}.
\end{equation}
Let $\tau_{\calN}$ be the $\rho_{\calN}$-quantile of $T_{\calN}$ under the node
population.  The split is
\begin{align}
  C_L&=\{x\in C_{\calN}:q_{\calN}^{\top}\tilde x\le \tau_{\calN}\},\\
  C_R&=\{x\in C_{\calN}:q_{\calN}^{\top}\tilde x>\tau_{\calN}\}.
\end{align}
Quantiles control mass.  They do not decide whether the split is geometrically
meaningful.

\subsection{Geometry layer: choosing split directions}

\paragraph{GA-SPR-Mass.}
The mass-only baseline chooses fixed basis directions or seeded random axes.
This tests the equal-mass construction without claiming feature meaning.

\paragraph{GA-SPR-Local.}
The local variant chooses $q$ using unlabeled feature geometry.  Given candidate
directions $q^{(m)}$, each with its quantile split, choose the direction that
minimizes weighted within-cell variance
\begin{equation}
  \calL_{\mathrm{local}}(q)=\sum_{b\in\{L,R\}}\Prob(C_b\mid C_{\calN})
  \E[\norm{\tilde X-\mu_b}_2^2\mid \tilde X\in C_b].
\end{equation}
A margin term can be added to avoid thresholds that cut through high-density
regions.

\paragraph{GA-SPR-Expert.}
Expert weights can guide the split.  For an input-side expert matrix $W_e$,
define its principal spectral line by
\begin{equation}
  v_e\in\arg\max_{\norm{v}_2=1}\norm{v^{\top}W_e}_2^2,\qquad
  P_e=v_ev_e^{\top}.
\end{equation}
The projector $P_e$ removes the arbitrary sign of the singular vector.  Expert
projectors can be grouped by a sign-invariant distance such as
\begin{equation}
  d(e,j)=1-\mathrm{tr}(P_eP_j)=1-(v_e^{\top}v_j)^2.
\end{equation}
After grouping experts into $\calE_L$ and $\calE_R$, GA-SPR-Expert chooses a
feature split direction by either using an oriented router/MPI row when available
or by selecting a locality direction whose resulting cells align with the two
expert subspaces.  In the simplest projective form, a candidate direction can be
extracted from the leading eigenspace of
\begin{equation}
  \Delta_{\calN}=\frac{1}{|\calE_L|}\sum_{e\in\calE_L}P_e
  -\frac{1}{|\calE_R|}\sum_{e\in\calE_R}P_e.
\end{equation}
The sign of this direction is immaterial for equal mass; it only determines
which child is named left.  Quantile thresholds still determine the child masses.

\subsection{Preservation layer: margin-aware updates}

The route margin of a tree router is
\begin{equation}
  m_{\mathrm{tree}}(x)=\min_{\calN\in\mathrm{path}(x)}
  |q_{\calN}^{\top}\tilde x-\tau_{\calN}|.
\end{equation}
Naive recalibration may move thresholds across many samples and cause route
flips.  GA-SPR therefore treats updates as controlled interventions.  A threshold
or direction update is accepted only if the estimated route-change budget on a
held-out calibration buffer is below a chosen tolerance.  A hysteresis band can
keep points near old boundaries fixed unless the new route is confident.  These
rules are not part of the equal-mass proof; they are a training-preservation
mechanism to be tested experimentally.

\section{Theory}

\subsection{Equal mass from recursive quantiles}

\begin{theorem}[Recursive quantile dispatch]
\label{thm:recursive_quantile}
Assume each internal node $\calN$ chooses a threshold satisfying
\begin{equation}
  \Prob(C_L\mid C_{\calN})=\frac{|\calE_L|}{|\calE_{\calN}|},\qquad
  \Prob(C_R\mid C_{\calN})=\frac{|\calE_R|}{|\calE_{\calN}|}.
\end{equation}
If the root contains all $E$ experts and every leaf contains one expert, then
each expert cell has mass $1/E$.
\end{theorem}

\begin{proof}
For any node $\calN$, prove by induction that
$\Prob(C_{\calN})=|\calE_{\calN}|/E$.  The root satisfies this.  If the parent
satisfies it, then the left child satisfies
\begin{equation}
\Prob(C_L)=\Prob(C_{\calN})\Prob(C_L\mid C_{\calN})
=\frac{|\calE_{\calN}|}{E}\frac{|\calE_L|}{|\calE_{\calN}|}
=\frac{|\calE_L|}{E}.
\end{equation}
The right child is identical.  A leaf has one expert, so its mass is $1/E$.
\end{proof}

\subsection{Empirical quantiles and atom boundary}

The theorem is exact for population quantiles.  With a finite calibration set,
empirical CDF concentration gives the usual $O(n^{-1/2})$ split error for a
continuous scalar projection \citep{dvoretzky1956asymptotic}.  For discrete
features, exact deterministic splitting may be impossible.

\begin{proposition}[Atom impossibility]
Suppose an atom $x_0$ has probability $\Prob(X=x_0)>1/E$.  Any deterministic
feature-only router that sends identical features to the same expert cannot have
$\Prob(g(X)=e)=1/E$ for all $e$.
\end{proposition}

\begin{proof}
All mass at $x_0$ must be sent to a single expert $e_0$.  Hence
$\Prob(g(X)=e_0)\ge \Prob(X=x_0)>1/E$, contradicting equal mass.
\end{proof}

Thus discrete experiments require explicit tie handling or an approximate target.
Feature-derived lexicographic keys are allowed; sample-index tie breaking is not
a deterministic feature router.

\subsection{Coherence under projected separation}

Equal mass does not imply feature coherence.  The next theorem states the mild
condition under which a split is coherent.

\begin{assumption}[Projected separation]
At node $\calN$, suppose the feature population contains two latent groups
$G\in\{L,R\}$ that should be assigned to the two children, with
$\Prob(G=L\mid C_{\calN})=\rho_{\calN}$.  A direction $q$ and threshold $\tau$
have projected separation error $\varepsilon$ if
\begin{equation}
  \Prob(q^{\top}\tilde X>\tau\mid G=L,C_{\calN})\le \varepsilon,\qquad
  \Prob(q^{\top}\tilde X\le\tau\mid G=R,C_{\calN})\le \varepsilon.
\end{equation}
\end{assumption}

\begin{proposition}[One-node coherence]
Under projected separation error $\varepsilon$, the node split disagrees with the
latent binary grouping with probability at most $\varepsilon$.
\end{proposition}

\begin{proof}
The disagreement event is the union of the left-group tail above the threshold
and the right-group tail below the threshold, weighted by the group priors.  Each
conditional error is at most $\varepsilon$, so the total error is at most
$\varepsilon$.
\end{proof}

For a tree, the route-to-feature error is at most the sum of separation errors
along the path.  This is the theoretical reason that good dispatch needs more
than isotropy: it needs directions whose projected splits separate the relevant
feature groups.

\subsection{Expert alignment as a sufficient source of split directions}

The expert-guided variant is justified only when expert weights contain useful
geometry.  The following perturbation statement makes this explicit.

\begin{assumption}[Aligned expert split]
For a node split, let $q^\star$ be an ideal feature-separating direction with
threshold $\tau^\star$.  A candidate direction $q$ from expert weights is
$\delta$-aligned if $\norm{q-q^\star}_2\le\delta$.  Features are bounded by
$\norm{\tilde X}_2\le R$ with probability one, and have margin
\begin{equation}
  m^\star(X)=|{q^\star}^{\top}\tilde X-\tau^\star|.
\end{equation}
\end{assumption}

\begin{proposition}[Expert-direction perturbation]
If $\norm{q-q^\star}_2\le\delta$ and the threshold is unchanged, then the split
can disagree with the ideal split only when $m^\star(X)\le \delta R$.  Therefore
\begin{equation}
  \Prob(g_q(X)\ne g_{q^\star}(X))
  \le \Prob(m^\star(X)\le \delta R).
\end{equation}
\end{proposition}

\begin{proof}
For any $X$, the projection perturbation is bounded by
$|(q-q^\star)^{\top}\tilde X|\le \delta R$.  If the ideal margin is larger than
this perturbation, the sign relative to the threshold cannot change.
\end{proof}

This proposition explains the role of expert weights.  They are useful when they
produce split directions close to feature-separating directions and when the
feature distribution has low boundary mass.

\subsection{Preservation under bounded drift}

\begin{theorem}[Margin preservation]
Let $g_0$ be a calibrated tree router with unit split directions and thresholds
$\tau_{\calN}$.  Let $X_t=X_0+\Delta$ be a drifted feature.  If
\begin{equation}
  \norm{\Delta}_2 < m_{\mathrm{tree}}(X_0),
\end{equation}
then $g_t(X_t)=g_0(X_0)$ under the same tree.
\end{theorem}

\begin{proof}
At each node on the original path, the split score changes by at most
$|q_{\calN}^{\top}\Delta|\le\norm{\Delta}_2$.  If this is smaller than the
minimum distance to all thresholds on the path, no binary decision changes.
\end{proof}

Consequently, if $\norm{\Delta}_2\le r$ with high probability, the route flip
probability is bounded by $\Prob(m_{\mathrm{tree}}(X_0)\le r)$ plus the drift
failure probability.  Training preservation is therefore a margin-and-drift
problem, not a consequence of equal-mass initialization.

\section{Experiments}
\label{sec:experiments}

This draft intentionally leaves the experimental results empty.  The protocols
below define what must be run before claiming empirical validation.

\subsection{Static equal-mass dispatch}

\paragraph{Goal.} Verify that recursive quantile routing constructs equal-mass
cells on calibration data and generalizes to holdout data.

\paragraph{Data.} Gaussian, spherical, and whitened feature distributions.
Report both continuous-density cases and discrete/atom-supported boundary cases.

\paragraph{Methods.} Random linear router, regular-simplex prototype router,
SPR-Mass, GA-SPR-Local, GA-SPR-Expert, and balanced assignment control.

\paragraph{Metrics.} Calibration LoadCV, holdout LoadCV, maximum usage error,
calibration-to-holdout gap, tie rate, atom mass, and margin p05.

\subsection{Assumption ladder for feature meaning}

\paragraph{Goal.} Determine which additional assumptions are needed beyond
spectral isotropy for meaningful feature partition.

\paragraph{Data regimes.}
\begin{enumerate}
\item Isotropic but unstructured distributions.
\item Isotropic and clusterable mixtures $X=\mu_F+\epsilon$.
\item Expert-aligned mixtures where expert spectral prototypes correlate with
latent feature directions.
\item Low-margin mixtures where equal-mass dispatch is possible but unstable.
\end{enumerate}

\paragraph{Metrics.} LoadCV, route-feature NMI, purity, within-cell variance,
cell-center separation, and margin p05.

\subsection{Expert-weight-guided routing}

\paragraph{Goal.} Test whether expert weights provide useful geometry after
fixing sign ambiguity.

\paragraph{Required fixes before running.} Use sign-invariant projective
prototypes or a deterministic orientation rule. Report prototype spread using
nearest-neighbor cosine, maximum cosine, and effective prototype rank.

\paragraph{Methods.} Expert-SVD top-1, MPI top-1, SPR-Mass, GA-SPR-Local,
GA-SPR-Expert, and GA-SPR-Expert with local refinement.

\paragraph{Metrics.} LoadCV, prototype spread, cell-prototype alignment,
route-feature NMI, purity, and margin.

\subsection{Margin-preserved training}

\paragraph{Goal.} Test whether a calibrated partition is preserved during a
controlled hard-routing training task.

\paragraph{Task.} Synthetic latent feature classes, with class-dependent target
maps $Y=A_FX+\epsilon$.  Feature labels are hidden from router construction and
training, but used for evaluation.

\paragraph{Routers.} Random linear top-1, true trainable linear top-1 with
auxiliary load-balancing loss, balanced assignment control, Frozen GA-SPR,
GA-SPR-init-linear, expert-guided GA-SPR, and margin-aware update GA-SPR.

\paragraph{Metrics.} Validation loss, LoadCV, route-feature NMI, feature purity,
route flip rate, margin p05, representation drift, and expert specialization.

\subsection{Ablations}

Ablate calibration size, expert count, dimension, feature atom mass, split policy,
expert-weight quality, margin-update budget, hysteresis width, and drift strength.

\section{Discussion}

The central lesson is that uniform feature dispatch has layers.  Equal mass is a
solvable construction problem: recursive quantiles do it under standard
regularity conditions.  Meaningful partition is a structure problem: isotropy
alone cannot tell the router which features belong together.  Preservation is a
dynamics problem: even a good initial partition can be lost if drift exceeds the
routing margin.

The next useful assumption is therefore not full spherical uniformity, which is
strong, nor covariance isotropy alone, which is too weak.  A productive middle
regime is globally isotropic but locally structured features, possibly with
expert weights that carry an aligned geometric prior.  This regime is close to
what controlled slot-feature experiments try to capture: identify the
route-relevant feature population, build equal-mass cells, attach cells to
experts, and then check whether training preserves the assignment.

\section{Scope and Limitations}

GA-SPR is not a claim that equal-mass routing discovers semantic experts in
general.  Without clusterability, expert alignment, or another structure prior,
there may be no meaningful feature partition to recover.  The theory also assumes
population or well-estimated quantiles; finite calibration, atoms, and
distribution shift must be audited.  Finally, margin-aware updates are proposed
as a training mechanism, not proven to solve arbitrary representation drift.

\section{Conclusion}

We refine the uniform-routing problem into a precise next-stage target:
construct routers that achieve equal mass, feature coherence, and training
preservation in a spectrally normalized feature space.  Quantiles solve the mass
part; split directions solve the geometry part only under additional structure;
and margins determine preservation.  GA-SPR organizes these roles into a single
router family and gives the theory and experiment plan needed to test the next
iteration.

\bibliographystyle{iclr2026_conference}
\bibliography{references}

\end{document}
